\documentclass[12pt]{article}
\usepackage[utf8]{inputenc}
\usepackage[spanish, es-tabla]{babel} 
\usepackage{graphicx}
\usepackage{amsmath}
\usepackage{amssymb}
\usepackage{geometry}
\usepackage{hyperref}
\usepackage{booktabs}
\usepackage{array}
\usepackage{xcolor}
\usepackage{float}
\usepackage{setspace}
\usepackage{caption}
\usepackage{subcaption}
\usepackage{longtable}
\usepackage{multirow}
\usepackage{makecell}

% Configuración de márgenes
\geometry{a4paper, margin=1in}

% Configuración de hipervínculos
\hypersetup{
    colorlinks=true,
    linkcolor=blue,
    filecolor=magenta,      
    urlcolor=cyan,
    pdftitle={Análisis Estadístico NBA},
    pdfpagemode=FullScreen,
}

% Comando personalizado para preguntas
\newcommand{\pregunta}[1]{\textbf{\large #1}}

% Configuración de espaciado
\onehalfspacing

\begin{document}

% --- Portada ---
\begin{titlepage}
    \centering
    \vspace*{1cm}
    \LARGE\textbf{Análisis Estadístico de la NBA: Reporte Completo}\\
    \vspace{0.5cm}
    \large\textbf{Análisis Exploratorio de Datos y Estadísticos Descriptivos}\\
    \vspace{1.5cm}
    \normalsize Dataset: Estadísticas Temporada Regular NBA (2010-2024)\\
    \vspace{2cm}
    \textbf{Autores:}\\
    Daniela De La Caridad Guerrero Álvarez (C311)\\
    Sammy Raul Sosa Justiz (C312)\\
    Rubén Martínez Rojas (C311)\\
    \vfill
    \today
\end{titlepage}

% --- Índice ---
\tableofcontents
\newpage

% --- Introducción ---
\section*{Introducción}
\addcontentsline{toc}{section}{Introducción}
Este dataset permite analizar tendencias históricas, patrones de rendimiento y factores asociados al éxito en la NBA. Las preguntas propuestas abordan temas centrales del baloncesto moderno: la importancia del triple, la evolución ofensiva y los impulsores estadísticos del rendimiento en partidos.

% --- Preguntas de Investigación ---
\subsection*{\pregunta{Pregunta 1}}

\textbf{¿Existe una correlación significativa entre el número de triples anotados (FG3M) y la victoria (WL) de un equipo en un partido, controlando por la efectividad general de tiro (FG\_PCT)?}

\vspace{0.3cm}

\textbf{Relevancia:} El baloncesto moderno ha evolucionado hacia un mayor énfasis en el tiro de tres puntos. Entender si un alto volumen de triples está asociado con mayores probabilidades de ganar, independientemente de la efectividad general, es clave para estrategias de juego y análisis de desempeño. Esto ayuda a evaluar si la filosofía de ``vivir o morir por el triple'' es estadísticamente sólida.

\vspace{0.3cm}

\textbf{Variables involucradas:}
\begin{itemize}
    \setlength\itemsep{0.1em}
    \item \textbf{Variable respuesta:} WL (Win/Loss)
    \item \textbf{Variables predictoras:} FG3M (triples anotados), FG\_PCT (porcentaje total de tiros de campo)
\end{itemize}

\vspace{0.3cm}

\textbf{Técnicas posibles:}
\begin{itemize}
    \setlength\itemsep{0.1em}
    \item Regresión logística
    \item Análisis de correlación parcial
    \item Pruebas de chi-cuadrado
\end{itemize}

% Pregunta 2
\subsection*{\pregunta{Pregunta 2}}

\textbf{¿Cómo ha cambiado la media de puntos por partido (PTS) y el porcentaje de tiros de tres puntos intentados (FG3A/FGA) a lo largo de las temporadas (SEASON\_YEAR), y existe una relación temporal significativa entre ambas variables?}

\vspace{0.3cm}

\textbf{Relevancia:} La NBA ha experimentado una ``revolución del triple'' en la última década. Analizar la tendencia temporal en puntos y en el uso del triple permite cuantificar la transformación ofensiva de la liga. Esto es útil para entrenadores, analistas y medios para contextualizar récords ofensivos y evoluciones estratégicas.

\vspace{0.3cm}

\textbf{Variables involucradas:}
\begin{itemize}
    \setlength\itemsep{0.1em}
    \item \textbf{Variable temporal:} SEASON\_YEAR
    \item \textbf{Variables de interés:} PTS (puntos), FG3A (intentos de triple), FGA (intentos totales)
\end{itemize}

\vspace{0.3cm}

\textbf{Técnicas posibles:}
\begin{itemize}
    \setlength\itemsep{0.1em}
    \item Series temporales
    \item Regresión lineal
    \item Análisis de tendencia
    \item Comparación de medias por temporada (ANOVA)
\end{itemize}

% Pregunta 3
\subsection*{\pregunta{Pregunta 3}}

\textbf{¿Qué combinación de estadísticas de juego (por ejemplo, asistencias AST, rebotes REB, robos STL, porcentaje de tiros FG\_PCT) predice de manera más robusta la diferencia de puntos (PLUS\_MINUS) en un partido?}

\vspace{0.3cm}

\textbf{Relevancia:} El plus-minus refleja el impacto neto de un equipo durante el partido. Identificar qué estadísticas de box score contribuyen más a un plus-minus positivo ayuda a priorizar métricas de rendimiento y diseñar estrategias de equipo. Es fundamental para el análisis avanzado y la evaluación de jugadores.

\vspace{0.3cm}

\textbf{Variables involucradas:}
\begin{itemize}
    \setlength\itemsep{0.1em}
    \item \textbf{Variable respuesta:} PLUS\_MINUS
    \item \textbf{Variables predictoras:} AST, REB, STL, BLK, TOV, FG\_PCT, FT\_PCT, FG3\_PCT, etc.
\end{itemize}

\vspace{0.3cm}

\textbf{Técnicas posibles:}
\begin{itemize}
    \setlength\itemsep{0.1em}
    \item Regresión lineal múltiple
    \item Selección de modelos (stepwise, lasso)
    \item Análisis de componentes principales
\end{itemize}

% Sección: Metodología Propuesta
\section*{Metodología Propuesta}

\subsection*{Fuente de Datos}
\begin{itemize}
    \item \textbf{Dataset:} regular\_season\_totals\_2010\_2024.csv
    \item \textbf{Período:} Temporadas 2010-2024 de la NBA
    \item \textbf{Observaciones:} Estadísticas por partido a nivel de equipo
    \item \textbf{Variables:} 52 columnas incluyendo estadísticas básicas, avanzadas y rankings
\end{itemize}

% --- Información General ---
\section{Información General del Dataset}
\begin{table}[H]
\centering
\begin{tabular}{@{}ll@{}}
\toprule
\textbf{Característica} & \textbf{Valor} \\ \midrule
Total de registros & 1,231 \\
Total de variables & 52 \\
Período cubierto & 2010-2024 \\
Equipos representados & 30 equipos de la NBA \\
Valores nulos totales & 202 (0.31\%) \\ \bottomrule
\end{tabular}
\caption{Características generales del dataset}
\end{table}

% --- Clasificación ---
\section{Clasificación de Variables}

\subsection{Variables Cualitativas}
\begin{itemize}
\item \textbf{Identificación}: SEASON\_YEAR, TEAM\_ID, TEAM\_ABBREVIATION, TEAM\_NAME
\item \textbf{Juego}: GAME\_ID, GAME\_DATE, MATCHUP, WL
\item \textbf{Control}: AVAILABLE\_FLAG
\item \textbf{Rankings categóricos}: GP\_RANK, W\_RANK, L\_RANK
\end{itemize}

\subsection{Variables Cuantitativas Discretas}
\begin{itemize}
\item \textbf{Conteos estadísticos}: FGM, FGA, FG3M, FG3A, FTM, FTA
\item \textbf{Rebotes}: OREB, DREB
\item \textbf{Otras estadísticas}: AST, TOV, STL, BLK, BLKA, PF, PFD, PTS
\item \textbf{Variables de ranking discretas}: W\_PCT\_RANK, MIN\_RANK, FG\_PCT\_RANK, etc.
\end{itemize}

\subsection{Variables Cuantitativas Continuas}
\begin{itemize}
\item \textbf{Porcentajes}: FG\_PCT, FG3\_PCT, FT\_PCT
\item \textbf{Estadísticas continuas}: MIN, PLUS\_MINUS
\item \textbf{Estadísticas calculadas}: REB (suma de OREB + DREB)
\end{itemize}


% --- Estadísticos Descriptivos ---
\section{Estadísticos Descriptivos}

\subsection{Medidas de Tendencia Central y Dispersión}

\begin{table}[H]
\centering
\small
\begin{tabular}{|l|c|c|c|c|c|}
\hline
\textbf{Estadístico} & \textbf{PTS} & \textbf{REB} & \textbf{AST} & \textbf{FG\_PCT} & \textbf{PLUS\_MINUS} \\
\hline
Media & 124.38 & 43.33 & 28.95 & 0.500 & 17.47 \\
Mediana & 125.00 & 43.00 & 29.00 & 0.505 & 15.00 \\
Moda & 120.00 & 42.00 & 30.00 & 0.500 & 5.00 \\
Desv. Estándar & 15.25 & 7.43 & 6.06 & 0.046 & 17.31 \\
Varianza & 232.64 & 55.24 & 36.72 & 0.002 & 299.76 \\
Coef. Variación & 12.26\% & 17.15\% & 20.94\% & 9.20\% & 99.09\% \\
Mínimo & 89.00 & 26.00 & 6.00 & 0.357 & -28.00 \\
Máximo & 175.00 & 65.00 & 62.00 & 0.644 & 61.00 \\
Rango & 86.00 & 39.00 & 56.00 & 0.287 & 89.00 \\
Q1 & 115.00 & 39.00 & 25.00 & 0.473 & 5.00 \\
Q3 & 136.00 & 48.00 & 33.00 & 0.532 & 30.00 \\
IQR & 21.00 & 9.00 & 8.00 & 0.059 & 25.00 \\
\hline
\end{tabular}
\caption{Estadísticos descriptivos principales}
\end{table}

% Sección: Distribución de Variables Cualitativas
\section{Distribución de Variables Cualitativas}

\subsection{Resultados de Partidos (WL)}
\begin{table}[H]
\centering
\begin{tabular}{|c|c|c|}
\hline
\textbf{Resultado} & \textbf{Frecuencia} & \textbf{Porcentaje} \\
\hline
Victoria (W) & 615 & 49.96\% \\
Derrota (L) & 616 & 50.04\% \\
\hline
\textbf{Total} & 1,231 & 100.00\% \\
\hline
\end{tabular}
\caption{Distribución de resultados de partidos}
\end{table}

\subsection{Equipos con Más Registros}
\begin{table}[H]
\centering
\small
\begin{tabular}{|c|l|c|c|}
\hline
\textbf{Rank} & \textbf{Equipo} & \textbf{Registros} & \textbf{Porcentaje} \\
\hline
1 & HOU (Houston Rockets) & 92 & 7.47\% \\
2 & GSW (Golden State Warriors) & 90 & 7.31\% \\
3 & BOS (Boston Celtics) & 52 & 4.22\% \\
4 & POR (Portland Trail Blazers) & 36 & 2.92\% \\
5 & MIL (Milwaukee Bucks) & 35 & 2.84\% \\
6 & UTA (Utah Jazz) & 33 & 2.68\% \\
7 & DEN (Denver Nuggets) & 27 & 2.19\% \\
8 & LAC (LA Clippers) & 25 & 2.03\% \\
9 & NYK (New York Knicks) & 23 & 1.87\% \\
10 & ATL (Atlanta Hawks) & 23 & 1.87\% \\
\hline
\end{tabular}
\caption{Top 10 equipos con más registros en el dataset}
\end{table}


% --- TDEF ---
\section{Tablas de Distribución Empírica de Frecuencia (TDEF)}

\subsection{Distribución de PUNTOS (PTS)}
\begin{table}[H]
\centering
\footnotesize
\begin{tabular}{@{}lcccc@{}}
\toprule
\textbf{Intervalo} & \textbf{Frecuencia} & \textbf{Frec. \%} & \textbf{Frec. Acum.} & \textbf{Frec. Acum. \%} \\ \midrule
$[89.00, 99.75)$ & 35 & 2.84\% & 35 & 2.84\% \\
$[99.75, 110.50)$ & 138 & 11.21\% & 173 & 14.05\% \\
$[110.50, 121.25)$ & 335 & 27.21\% & 508 & 41.26\% \\
$[121.25, 132.00)$ & 360 & 29.24\% & 868 & 70.51\% \\
$[132.00, 142.75)$ & 213 & 17.30\% & 1,081 & 87.81\% \\
$[142.75, 153.50)$ & 106 & 8.61\% & 1,187 & 96.42\% \\
$[153.50, 164.25)$ & 35 & 2.84\% & 1,222 & 99.26\% \\
$[164.25, 175.00]$ & 9 & 0.73\% & 1,231 & 100.00\% \\ \bottomrule
\end{tabular}
\caption{TDEF para PTS (n=1,231, Rango=86, Amplitud=10.75)}
\end{table}

\subsection{TDEF para PORCENTAJE DE TIROS (FG\_PCT)}
\small
\begin{table}[H]
\centering
\small
\begin{tabular}{@{}lcccc@{}}
\toprule    
\textbf{Intervalo} & \textbf{ni} & \textbf{ni (\%)} & \textbf{Ni} & \textbf{Ni (\%)} \\ \midrule
$[0.357, 0.393)$ & 39 & 3.17\% & 39 & 3.17\% \\
$[0.393, 0.429)$ & 99 & 8.04\% & 138 & 11.21\% \\
$[0.429, 0.465)$ & 226 & 18.36\% & 364 & 29.57\% \\
$[0.465, 0.501)$ & 347 & 28.19\% & 711 & 57.76\% \\
$[0.501, 0.537)$ & 293 & 23.80\% & 1,004 & 81.56\% \\
$[0.537, 0.573)$ & 157 & 12.75\% & 1,161 & 94.31\% \\
$[0.573, 0.609)$ & 57 & 4.63\% & 1,218 & 98.94\% \\
$[0.609, 0.644]$ & 13 & 1.06\% & 1,231 & 100.00\% \\ \bottomrule
\end{tabular}
\caption{TDEF para FG\_PCT}
\end{table}

% Sección: Análisis de Correlaciones
\section{Análisis de Correlaciones}

\subsection{Matriz de Correlación}

\begin{table}[H]
\centering
\begin{tabular}{|l|c|c|c|c|c|}
\hline
 & \textbf{PTS} & \textbf{REB} & \textbf{AST} & \textbf{FG\_PCT} & \textbf{PLUS\_MINUS} \\
\hline
\textbf{PTS} & 1.0000 & 0.4523 & 0.5728 & 0.6731 & 0.6142 \\
\textbf{REB} & 0.4523 & 1.0000 & 0.3245 & 0.2104 & 0.3251 \\
\textbf{AST} & 0.5728 & 0.3245 & 1.0000 & 0.3427 & 0.4503 \\
\textbf{FG\_PCT} & 0.6731 & 0.2104 & 0.3427 & 1.0000 & 0.5214 \\
\textbf{PLUS\_MINUS} & 0.6142 & 0.3251 & 0.4503 & 0.5214 & 1.0000 \\
\hline
\end{tabular}
\caption{Matriz de correlación entre variables principales}
\end{table}

\subsection{Correlaciones Significativas}

\begin{itemize}
\item \textbf{PTS - FG\_PCT}: 0.6731 (correlación alta positiva) - La efectividad de tiro explica gran parte de la variación en puntos.

\item \textbf{PTS - AST}: 0.5728 (correlación moderada-alta positiva) - Las asistencias contribuyen significativamente a la generación de puntos.

\item \textbf{PTS - REB}: 0.4523 (correlación moderada positiva) - Los rebotes tienen menor correlación con puntos que las asistencias.

\item \textbf{FG\_PCT - PLUS\_MINUS}: 0.5214 (correlación moderada) - La efectividad de tiro está relacionada con el diferencial de puntos.
\end{itemize}

\subsection{Conclusiones}
El análisis demuestra que la NBA actual presenta una alta eficiencia de tiro ($\bar{x} = 50\%$). La fuerte correlación entre \texttt{FG\_PCT} y \texttt{PTS} ($r = 0.67$) confirma que la efectividad es el predictor primario del volumen anotador, por encima del volumen de rebotes.

% Sección: Análisis de Outliers
\section{Análisis de Outliers (Método IQR)}

\begin{table}[H]
\centering
\begin{tabular}{|l|c|c|c|c|c|}
\hline
\textbf{Variable} & \textbf{Total} & \textbf{Outliers} & \textbf{\%} & \textbf{Límite Inf.} & \textbf{Límite Sup.} \\
\hline
PTS & 1,231 & 47 & 3.82\% & 83.5 & 167.5 \\
REB & 1,231 & 36 & 2.92\% & 25.5 & 61.5 \\
AST & 1,231 & 25 & 2.03\% & 13.0 & 45.0 \\
PLUS\_MINUS & 1,231 & 83 & 6.74\% & -32.5 & 67.5 \\
\hline
\end{tabular}
\caption{Detección de outliers usando método IQR (Q1 - 1.5*IQR, Q3 + 1.5*IQR)}
\end{table}

\subsection{Partidos Extremos Identificados}
\begin{itemize}
\item \textbf{Máximo de puntos}: 175 puntos (outlier superior)
\item \textbf{Mínimo de puntos}: 89 puntos (no es outlier)
\item \textbf{Máximo de rebotes}: 65 rebotes (outlier superior)
\item \textbf{Máximo de asistencias}: 62 asistencias (outlier superior)
\item \textbf{Máximo plus/minus}: +61 (outlier superior)
\item \textbf{Mínimo plus/minus}: -28 (no es outlier)
\end{itemize}

\subsection{Análisis de Outliers (Método IQR)}
Se detectaron valores atípicos significativos en la variable \texttt{PLUS\_MINUS} (6.74\%), lo que refleja partidos con diferencias de puntuación inusualmente altas para el promedio de la liga.

% ====================================================
% ANÁLISIS ESTADÍSTICO - GENERADO DINÁMICAMENTE
% Ejecutar: python generar_todo.py antes de compilar
% ====================================================

\section{Análisis Estadístico de Preguntas de Investigación}

% AUTO-GENERADO por pregunta1_logistica.py
\subsection{Pregunta 1: Correlación FG3M vs Victoria}

\subsubsection{Análisis Chi-Cuadrado}

\textbf{Objetivo:} Determinar si existe asociación significativa entre FG3M y WL.

\begin{table}[H]
\centering
\begin{tabular}{|l|c|c|c|}
\hline
\textbf{FG3M} & \textbf{Derrota} & \textbf{Victoria} & \textbf{Total} \\
\hline
Bajo & 9633 & 6897 & 16530 \\
Medio & 6917 & 9238 & 16155 \\
Alto & 108 & 523 & 631 \\
\hline
\end{tabular}
\caption{Tabla de contingencia FG3M vs WL}
\end{table}

\begin{table}[H]
\centering
\begin{tabular}{|l|c|}
\hline
\textbf{Estadístico} & \textbf{Valor} \\
\hline
$\chi^2$ & 1059.25 \\
Grados de libertad & 2 \\
Valor p & $<$0.001 \\
V de Cramér & 0.178 \\
\hline
\end{tabular}
\caption{Resultados Chi-Cuadrado ($\alpha$ = 0.05)}
\end{table}

\subsubsection{Regresión Logística}

\textbf{Modelo:} $\log\left(\frac{P(Win)}{1-P(Win)}\right) = \beta_0 + \beta_1 \cdot FG3M + \beta_2 \cdot FG\_PCT$

\begin{table}[H]
\centering
\begin{tabular}{|l|c|c|}
\hline
\textbf{Variable} & \textbf{Coeficiente} & \textbf{Odds Ratio} \\
\hline
FG3M & 0.21 & 1.23 \\
FG\_PCT & 1.08 & 2.93 \\
\hline
\end{tabular}
\caption{Coeficientes de regresión logística}
\end{table}

\textbf{Conclusión:} La filosofía de vivir o morir por el triple tiene respaldo estadístico.

% AUTO-GENERADO por pregunta2_analisis.py
\subsection{Pregunta 2: Evolución Temporal de PTS y Triples}

\subsubsection{Regresión Lineal Simple}

\textbf{Objetivo:} Cuantificar la relación temporal entre la temporada y PTS/FG3A.

\begin{table}[H]
\centering
\begin{tabular}{|l|c|c|c|c|}
\hline
\textbf{Variable} & \textbf{Pendiente} & \textbf{$R^2$} & \textbf{Valor p} & \textbf{Sig.} \\
\hline
PTS & +1.47 pts/temp & 0.922 & $<$0.001 & Sí \\
FG3A/FGA & +1.58\%/temp & 0.952 & $<$0.001 & Sí \\
\hline
\end{tabular}
\caption{Regresión lineal simple ($\alpha$ = 0.05)}
\end{table}

\subsubsection{ANOVA por Temporada}

\textbf{Objetivo:} Verificar diferencias significativas entre temporadas.

\begin{table}[H]
\centering
\begin{tabular}{|l|c|c|c|}
\hline
\textbf{Variable} & \textbf{F} & \textbf{Valor p} & \textbf{Decisión} \\
\hline
PTS & 647.86 & $<$0.001 & Rechazar $H_0$ \\
FG3A/FGA & 1950.55 & $<$0.001 & Rechazar $H_0$ \\
\hline
\end{tabular}
\caption{ANOVA por temporada ($\alpha$ = 0.05)}
\end{table}

\textbf{Conclusión:} La NBA ha experimentado una transformación ofensiva significativa, con aumentos de 1.5 puntos y 1.6\% en triples por temporada.

% AUTO-GENERADO por pregunta3_multiple.py
\subsection{Pregunta 3: Predictores de PLUS\_MINUS}

\subsubsection{Regresión Lineal Múltiple}

\textbf{Objetivo:} Identificar qué estadísticas predicen mejor PLUS\_MINUS.

\begin{table}[H]
\centering
\begin{tabular}{|l|c|c|c|c|}
\hline
\textbf{Variable} & \textbf{$\beta$} & \textbf{Error} & \textbf{t} & \textbf{p-valor} \\
\hline
AST & 0.809 & 0.040 & 20.17 & $<$0.001$^{***}$ \\
REB & 0.490 & 0.035 & 14.15 & $<$0.001$^{***}$ \\
STL & 1.483 & 0.116 & 12.75 & $<$0.001$^{***}$ \\
TOV & -1.562 & 0.080 & -19.62 & $<$0.001$^{***}$ \\
FG\_PCT & 85.228 & 4.715 & 18.08 & $<$0.001$^{***}$ \\
\hline
\end{tabular}
\caption{Coeficientes de regresión. $^{***}$p$<$0.001}
\end{table}

\begin{table}[H]
\centering
\begin{tabular}{|l|c|}
\hline
\textbf{Estadístico} & \textbf{Valor} \\
\hline
$R^2$ & 0.549 \\
$R^2$ ajustado & 0.547 \\
F-estadístico & 298.36 \\
\hline
\end{tabular}
\caption{Bondad de ajuste}
\end{table}

\textbf{Conclusión:} El modelo explica el 54.9\% de la variabilidad. FG\_PCT es el predictor más fuerte.






% --- Conclusiones Finales ---
\section*{Conclusiones Finales}
\addcontentsline{toc}{section}{Conclusiones Finales}

\begin{itemize}
\item El análisis exploratorio revela distribuciones normales para la mayoría de las variables estadísticas principales.
\item Las correlaciones más fuertes se observan entre puntos y porcentaje de tiro, confirmando la importancia de la eficiencia en el baloncesto moderno.
\item La detección de outliers indica partidos extremos que podrían merecer análisis más detallado.
\item El dataset presenta alta calidad de datos con valores nulos mínimos y consistencia en las relaciones entre variables.
\end{itemize}

\end{document}
